\documentclass[11pt]{article}
\usepackage{amsmath,amssymb,amsthm,hyperref}
\newtheorem{theorem}{Theorem}
\newtheorem{lemma}[theorem]{Lemma}
\title{Rational quadrature: a denominator-depth refinement}
\author{Companion to the general-prime lower-bound note}
\date{30 September 2026}
\begin{document}
\maketitle

Use the notation and the classical unit-indicator polynomial $H_{r,p}$
from \texttt{../rational\_designs/p\_adic\_lower\_bound.tex}.
The following refinement uses exactly the same composition argument.

\begin{lemma}
For a positive integer $N$ divisible by $p^a$, $a\ge1$, put
\[
 E_{r,p,a}=\left\lceil
       \frac{(1-p^{-a})(r+1)}{p-1}\right\rceil-1.
\]
Then
\[
 v_p\left(\int_0^1H_{r,p}(Nx)\,dx\right)\ge E_{r,p,a}.
\]
\end{lemma}
\begin{proof}
For a primitive $p$th root of unity $\zeta$, let
$\pi_h=\zeta^h-1$ and normalize $w(\pi_h)=1$, so
$w(q)=(p-1)v_p(q)$ for rational $q$. Set $\alpha=1-p^{-a}$.
Every positive-degree coefficient of
\[
 W_{N,h}(t)=(1+\pi_h t)^N-1
\]
satisfies
\[
 w([t^k]W_{N,h})\ge\lceil\alpha k\rceil+1.
\]
For $k\ge p^a$ the factor $\pi_h^k$ is enough, since
$\alpha k\le k-1$. For $1\le k<p^a$, the identity
$\binom Nk=(N/k)\binom{N-1}{k-1}$ shows that $p\mid\binom Nk$;
thus the valuation is at least $k+p-1\ge\lceil\alpha k\rceil+1$.

Write $G(z)=z/\log(1+z)=\sum_jg_jz^j$. The companion note proves
$w(g_j)\ge-j$. Consequently the coefficient $A_{k,h}$ in
\[
 \int_0^1(1+\pi_h t)^{Nx}\,dx
   =G(W_{N,h}(t))=\sum_kA_{k,h}t^k
\]
has valuation at least $\lceil\alpha k\rceil$.
The same uniform-convergence argument used in that note gives
$\sum_kA_{k,h}=1$ at $t=1$, because $p\mid N$.
Therefore
\[
 \int_0^1H_{r,p}(Nx)\,dx
   =\frac1p\sum_{h=1}^{p-1}\sum_{k>r}A_{k,h}
\]
has $w$-valuation at least $\lceil\alpha(r+1)\rceil-(p-1)$.
It is rational, so its integral-valued $p$-adic valuation is at least
$\lceil\alpha(r+1)/(p-1)\rceil-1$, as claimed.
\end{proof}

\begin{theorem}
Suppose rational nodes $x_1,\ldots,x_K$ give a degree-$r$ equal-weight
uniform interval rule. Fix a prime $p$, and suppose
\[
 a=\max_i\max(0,-v_p(x_i))\ge1.
\]
Then
\[
 K\ge p^{E_{r,p,a}}.
\]
More precisely, if $q$ is the number of indices $i$ with
$v_p(x_i)=-a$, then
\[
 p^{\min(\lfloor r/(p-1)\rfloor,\ v_p(K)+E_{r,p,a})}\mid q.
\]
\end{theorem}
\begin{proof}
The scaled nodes $y_i=p^a x_i$ are $p$-integral, and at least one is a
unit. Apply quadrature exactness to $H_{r,p}(p^a x)$ and the unit filter
modulo $p^{\lfloor r/(p-1)\rfloor}$. The lemma gives the divisibility.
Since $E_{r,p,a}\le\lfloor r/(p-1)\rfloor$ and $1\le q\le K$,
the displayed lower bound follows.
\end{proof}

For $a=1$, this recovers $E_{r,p,1}=\lfloor r/p\rfloor$.
For $p=2,a=2$, the exponent is
$\lceil3(r+1)/4\rceil-1$, asymptotic to $3r/4$.
For increasing denominator depth it approaches $r/(p-1)$.
The statement is conditional on the actual largest denominator valuation;
it does not improve the unconditional exponential lower bound because
that valuation may equal one.
\end{document}
